\documentclass[12pt]{article}
\usepackage{amsmath}
\usepackage{amsfonts}
\usepackage{float}
\usepackage{fancyhdr}
\usepackage{graphicx}

\usepackage{fontspec}
\usepackage{array}
\usepackage{longtable}
\usepackage{enumitem}
\setmainfont{DejaVu Serif}
\setsansfont{DejaVu Sans}
\setmonofont{DejaVu Sans Mono}
\graphicspath{{../../templates/}}
\setlength{\emergencystretch}{3em}
\setlist[enumerate]{label=\arabic*.,leftmargin=*}
\setlist[itemize]{leftmargin=*}
\usepackage{setspace}
\usepackage[colorlinks=true,linkcolor=blue, citecolor=red]{hyperref}
\usepackage{url}
\usepackage[top=.75in, left=.75in, right=.75in, bottom=1in]{geometry}


% For algorithm
\usepackage{algorithm}
\usepackage{algpseudocode}

% ============ CODE ============
\usepackage{listings}
\usepackage{xcolor}
\definecolor{codegreen}{rgb}{0,0.6,0}
\definecolor{codegray}{rgb}{0.5,0.5,0.5}
\definecolor{codepurple}{rgb}{0.58,0,0.82}
\definecolor{backcolour}{rgb}{0.95,0.95,0.92}

% Styling for the code.
\lstdefinestyle{mystyle}{
    backgroundcolor=\color{backcolour},   
    commentstyle=\color{codegreen},
    keywordstyle=\color{magenta},
    numberstyle=\tiny\color{codegray},
    stringstyle=\color{codepurple},
    basicstyle=\ttfamily\footnotesize,
    breakatwhitespace=false,         
    breaklines=true,                 
    captionpos=b,                    
    keepspaces=true,                 
    numbers=left,                    
    numbersep=5pt,                  
    showspaces=false,                
    showstringspaces=false,
    showtabs=false,                  
    tabsize=2
}
\lstset{style=mystyle}


% ============ HEADER AND FOOTER ============
% Header length
\setlength{\headheight}{29.43912pt}

% Footer page number would be on the lower-right corner
\pagestyle{fancy}
\fancyfoot{}
\fancyfoot[R]{Page \thepage}

% Disable indentation on new paragraphs
\setlength{\parindent}{0pt}

% Line spacing 1.5
\renewcommand{\baselinestretch}{1.5}

% Optional: graphic path
% \graphicspath{PATH_TO_GRAPHIC_FOLDER}

% To use Times font family, uncomment this row
% \usepackage{mathptmx}

% To use roman section / subsection, uncomment these rows
% \renewcommand{\thesection}{\Roman{section}}
% \renewcommand{\thesubsection}{\thesection.\Roman{subsection}}

% Define course name, report name and report title.
\newcommand{\coursename}{Introduction of Software Engineering}
\newcommand{\reportname}{Studitopia}
\newcommand{\reporttitle}{\href{https://github.com/RoundFloor/Studitopia.git}{Studitopia}}

\newcommand{\studentname}{%
Nguyen Thanh An (24125024)\\
Nguyen Chi Bao (24125025)\\
Tran Nguyen Khanh Binh (24125026)\\
Tran Canh Anh Tuan (24125048)\\
Nguyen Le Quoc Huy (24125100)
}
\newcommand{\teachername}{Ho Tuan Thanh, Nguyen Ngoc Minh Thu}

% Header
\lhead{\reporttitle}
\rhead{
University of Science - VNUHCM\\
\coursename
}

% Footer
\newcommand{\leftfooter}{}
\lfoot{\leftfooter}

% ============ DOCUMENT ============
\begin{document}

\pagenumbering{roman}
\begin{titlepage}
\begin{spacing}{1.0}
\newcommand{\HRule}{\rule{\linewidth}{0.5mm}}
\centering

\textsc{\LARGE VNU-HCM}\\[0.5cm]
\textsc{\Large University of Science}\\[0.5cm]
\textsc{\large Faculty of Information Technology}\\[0.5cm]
\textsc{Computer Science}\\[0.5cm]

\HRule \\[0.4cm]
{ 
\huge{\bfseries{\reporttitle}}\\[0.5cm]
\large{\bfseries{Project: \reportname}}
}\\[0.4cm]
\HRule \\[0.5cm]

\textbf{\large Subject: \coursename}\\[0.5cm]

\begin{minipage}[t]{0.55\textwidth}
\begin{flushleft} \large
\emph{Students:}\\
\studentname
\end{flushleft}
\end{minipage}
~
\begin{minipage}[t]{0.4\textwidth}
\begin{flushright} \large
\emph{Supervising Teachers:} \\
\teachername
\end{flushright}
\end{minipage}\\[1cm]

{\large \today}\\[1cm]

\includegraphics[scale=.20]{img/hcmus-logo.png}\\[1cm] 

\vfill
\end{spacing}
\end{titlepage}	

\tableofcontents
\pagebreak

\pagenumbering{arabic}
\setcounter{page}{1}

% ============ MARKDOWN CONTENT: PA2/software_development_plan.md ============
\hypertarget{sec:group-5---alien-force}{}
\section*{Group 5 - Alien Force}
\par\medskip
\hypertarget{sec:studitopia-development-plan-small-project}{}
\section*{Studitopia   Development Plan (Small Project)}
\par\medskip
\hypertarget{sec:version-11}{}
\subsection*{\textbf{Version:} \textless{}1.1\textgreater{}}
\par\medskip
\hypertarget{sec:revision-history}{}
\addcontentsline{toc}{section}{Revision History}
\section*{Revision History}
\par\medskip
\begingroup
\scriptsize
\renewcommand{\baselinestretch}{1.0}\selectfont
\renewcommand{\arraystretch}{1.15}
\setlength{\tabcolsep}{4pt}
\begin{longtable}{|>{\raggedright\arraybackslash}p{0.195\textwidth}|>{\raggedright\arraybackslash}p{0.156\textwidth}|>{\raggedright\arraybackslash}p{0.317\textwidth}|>{\raggedright\arraybackslash}p{0.250\textwidth}|}
\hline
Date & Version & Description & Author \\ \hline
\endhead
29/\allowbreak{}Oct/\allowbreak{}2026 & 1.0 & Initial draft & Tran Canh Anh Tuan \\ \hline
30/\allowbreak{}Oct/\allowbreak{}2026 & 1.1 & Plan the software development & Nguyen Thanh An \\ \hline
\end{longtable}
\endgroup
\par\medskip
\hypertarget{sec:1-introduction}{}
\addcontentsline{toc}{section}{1. Introduction}
\section*{1. Introduction}
\par\medskip
\textbf{Performed by:} [Tuan] \textbar{} \textbf{Reviewed by:} [An] \textbar{} \textbf{Edited by:} [An]
\par\medskip
This section provides an overview of the entire document. It includes the purpose, scope, definitions, acronyms, abbreviations, references, and overview of this plan.
\par\medskip
\hypertarget{sec:11-purpose}{}
\addcontentsline{toc}{subsection}{1.1 Purpose}
\subsection*{1.1 Purpose}
\par\medskip
The purpose of this document is to gather all information necessary to control the project. It describes the approach to the development of the software and direct the development effort of the group.
\par\medskip
The following people use the Software Development Plan:
\par\medskip
\begin{itemize}
\item The product owner, scrum master and project manager use it to plan the project schedule and resource needs, and to track progress against the schedule.
\par
\item Project team members use it to understand what they need to do, when they need to do it, and what other activities they are dependent upon.
\par
\end{itemize}
\hypertarget{sec:12-scope}{}
\addcontentsline{toc}{subsection}{1.2 Scope}
\subsection*{1.2 Scope}
\par\medskip
This document describes the overall plan to be used by the Studitopia project, including development, testing and deployment of the product which are interleaved in small iterations. The details of the individual iterations will be described in the Iteration Plans.
\par\medskip
The plans as outlined in this document are based upon the product requirements as defined in the Vision Document.
\par\medskip
\hypertarget{sec:13-overview}{}
\addcontentsline{toc}{subsection}{1.3 Overview}
\subsection*{1.3 Overview}
\par\medskip
The rest of this plan contains the following information:
\par\medskip
\begin{itemize}
\item \textbf{Project Overview} provides a description of the project's purpose, scope, and objectives. It also defines the deliverables that the project is expected to deliver.
\par
\item \textbf{Project Organization} describes the organizational structure of the project team.
\par
\item \textbf{Project Plan} contains the schedule and resources with their achievement criteria.
\par
\item \textbf{Project Monitoring and Control} describes how the project will be monitored and controlled.
\par
\end{itemize}
\hypertarget{sec:2-project-overview}{}
\addcontentsline{toc}{section}{2. Project Overview}
\section*{2. Project Overview}
\par\medskip
\textbf{Performed by:} [Tuan] \textbar{} \textbf{Reviewed by:} [An] \textbar{} \textbf{Edited by:} [An]
\par\medskip
\hypertarget{sec:21-project-purpose-scope-and-objectives}{}
\addcontentsline{toc}{subsection}{2.1 Project Purpose, Scope, and Objectives}
\subsection*{2.1 Project Purpose, Scope, and Objectives}
\par\medskip
Studitopia is a collaborative learning forum focused mainly on undergraduate students across several partner universities. The platform is organized as a forum with many sub-forums by major, and each major contains many courses. Users can join the forum to search for learning materials and past papers, and to discuss with other students about the learning process, study tips, and related topics. Users can share materials, host meetings within the forum, and track their personal progress against their university syllabus by recording credits and grades to see completed coursework and remaining requirements.
\par\medskip
Objectives:
\par\medskip
\begin{itemize}
\item Provide a course-centered forum where undergraduate students can access materials and past papers by major and course.
\par
\item Enable discussion and knowledge sharing among students about the learning process and study tips.
\par
\item Allow users to share learning materials and host meetings within the forum.
\par
\item Offer a progress tracking system tied to each university syllabus, where users input credits and grades to view their academic standing.
\par
\end{itemize}
\hypertarget{sec:22-assumptions-and-constraints}{}
\addcontentsline{toc}{subsection}{2.2 Assumptions and Constraints}
\subsection*{2.2 Assumptions and Constraints}
\par\medskip
Assumptions:
\par\medskip
\begin{itemize}
\item Users are undergraduate students enrolled at one of the partner universities.
\par
\item Users have access to a modern web browser and a stable internet connection.
\par
\item Staff members from partner universities may manage forums, update learning materials, moderate discussions, respond to students’ questions, and monitor user activity.
\par
\item The team will follow the Scrum process with 5 sprints which last about 12 weeks.
\par
\end{itemize}
Constraints:
\par\medskip
\begin{itemize}
\item Budget: student project with no dedicated funding; rely on free/\allowbreak{}open-source tools.
\par
\item Staff: 5 team members balancing coursework and project work.
\par
\item Schedule: project must be completed by the end of the semester (12 weeks).
\par
\item Platform: web browser only (no native mobile application in this phase).
\par
\item Technology: React + TypeScript + Vite (frontend), NestJS (backend), PostgreSQL (database only).
\par
\end{itemize}
\hypertarget{sec:23-project-deliverables}{}
\addcontentsline{toc}{subsection}{2.3 Project Deliverables}
\subsection*{2.3 Project Deliverables}
\par\medskip
\begingroup
\scriptsize
\renewcommand{\baselinestretch}{1.0}\selectfont
\renewcommand{\arraystretch}{1.15}
\setlength{\tabcolsep}{4pt}
\begin{longtable}{|>{\raggedright\arraybackslash}p{0.286\textwidth}|>{\raggedright\arraybackslash}p{0.407\textwidth}|>{\raggedright\arraybackslash}p{0.242\textwidth}|}
\hline
Deliverable & Description & Target \\ \hline
\endhead
Project proposal & The main ideas and features of the proposed application & Sprint 1 \\ \hline
Existing app survey & The survey about similar existing applications in the same domain & Sprint 1 \\ \hline
Vision Document & High-level needs, features, and non-functional requirements & Sprint 2 \\ \hline
Software Development Plan & Project plan covering organization, schedule, and control & Sprint 2 \\ \hline
Spec Kit Artifacts & Constitution and initialization files under /\allowbreak{}src & Sprint 2 \\ \hline
AI Usage Report \& Weekly Report & AI usage log and sprint meeting documentation & Each sprint \\ \hline
Use-case model and specification & User interaction with system object & Sprint 3 \\ \hline
Architecture diagram & The structure of the application in detail & Sprint 4 \\ \hline
Test cases and test plans & The detailed plan of testing process & Sprint 5 (final sprint) \\ \hline
Final Product Demo & Demonstration of all project features & Sprint 5 (final sprint) \\ \hline
\end{longtable}
\endgroup
\par\medskip
\hypertarget{sec:3-project-organization}{}
\addcontentsline{toc}{section}{3. Project Organization}
\section*{3. Project Organization}
\par\medskip
\textbf{Performed by:} [Bao] \textbar{} \textbf{Reviewed by:} [An] \textbar{} \textbf{Edited by:} []
\par\medskip
\hypertarget{sec:31-organizational-structure}{}
\addcontentsline{toc}{subsection}{3.1 Organizational Structure}
\subsection*{3.1 Organizational Structure}
\par\medskip
The project manager will create tasks in Jira and assigns to team members. There is also another member to review the work for each task. The review authorities for each tasks are also assigned by the project managers. All memebers work as full-stack developer, but each has specialist at one aspect, including front-end, tester, AI, UI design, database, and back-end. The member specialized one aspect will propose the overall structure and methods for implementing task relating to that aspect. Then, all members discuss to validate the proposed ideas.
\par\medskip
\hypertarget{sec:32-roles-and-responsibilities}{}
\addcontentsline{toc}{subsection}{3.2 Roles and Responsibilities}
\subsection*{3.2 Roles and Responsibilities}
\par\medskip
\begingroup
\scriptsize
\renewcommand{\baselinestretch}{1.0}\selectfont
\renewcommand{\arraystretch}{1.15}
\setlength{\tabcolsep}{4pt}
\begin{longtable}{|>{\raggedright\arraybackslash}p{0.221\textwidth}|>{\raggedright\arraybackslash}p{0.731\textwidth}|}
\hline
Person & Role \\ \hline
\endhead
Nguyễn Thanh An & Front-end Developer: Responsible for proposing methods to implement UI, designing UI components and their relations \\ \hline
Nguyễn Chí Bảo & Tester + AI Developer: Responsible for designing tests and architecture of AI server, researching suitable AI models suitable for available hardware, designing process to effectively and concurrently process data from multiple users using AI \\ \hline
Trần Nguyễn Khánh Bình & UI Designer: Responsible for designing the overall looking of UI including its components and pages alongside with its workflow \\ \hline
Trần Cảnh Anh Tuấn & Database Manager: Responsible for determining EER model, then establishing the relational schema while ensuring approriate normal form and usability \\ \hline
Nguyễn Lê Quốc Huy & Back-end Developer: Responsible for designing back-end API, communication with front-end, connection with database and AI server \\ \hline
\end{longtable}
\endgroup
\par\medskip
\hypertarget{sec:4-management-process}{}
\addcontentsline{toc}{section}{4. Management Process}
\section*{4. Management Process}
\par\medskip
\textbf{Performed by:} [An] \textbar{} \textbf{Reviewed by:} [Huy] \textbar{} \textbf{Edited by:} []
\par\medskip
\hypertarget{sec:41-project-estimates}{}
\addcontentsline{toc}{subsection}{4.1 Project Estimates}
\subsection*{4.1 Project Estimates}
\par\medskip
\hypertarget{sec:estimated-cost}{}
\addcontentsline{toc}{subsubsection}{Estimated cost}
\subsubsection*{Estimated cost}
\par\medskip
The project reuses existing infrastructure and open-source technologies, so the estimated additional development cost is \textbf{0 VND}.
\par\medskip
\begingroup
\scriptsize
\renewcommand{\baselinestretch}{1.0}\selectfont
\renewcommand{\arraystretch}{1.15}
\setlength{\tabcolsep}{4pt}
\begin{longtable}{|>{\raggedright\arraybackslash}p{0.320\textwidth}|>{\raggedright\arraybackslash}p{0.222\textwidth}|>{\raggedright\arraybackslash}p{0.394\textwidth}|}
\hline
Item & Estimated Cost & Explanation \\ \hline
\endhead
Server /\allowbreak{} Hosting & 0 VND & Our group owns server for hosting . \\ \hline
Domain & 0 VND & Domain has already been purchased. \\ \hline
File Storage & 0 VND & Existing storage is available in the server. \\ \hline
Database (PostgreSQL) & 0 VND & Open-source database. \\ \hline
Backend (NestJS) & 0 VND & Open-source framework. \\ \hline
Frontend (React + Vite) & 0 VND & Open-source technologies. \\ \hline
Real-time Communication & 0 VND & Implemented using WebSocket libraries. \\ \hline
Development and Testing Tools & 0 VND & Free or open-source tools will be used. \\ \hline
\end{longtable}
\endgroup
\par\medskip
Additional costs may arise if the existing server or storage cannot handle the expected workload, or if paid third-party services are required. These expenses will be estimated separately when necessary.
\par\medskip
\hypertarget{sec:estimated-schedule}{}
\addcontentsline{toc}{subsubsection}{Estimated schedule}
\subsubsection*{Estimated schedule}
\par\medskip
The project is expected to take \textbf{12 weeks}, divided into five two-week sprints.
\par\medskip
\begingroup
\scriptsize
\renewcommand{\baselinestretch}{1.0}\selectfont
\renewcommand{\arraystretch}{1.15}
\setlength{\tabcolsep}{4pt}
\begin{longtable}{|>{\raggedright\arraybackslash}p{0.170\textwidth}|>{\raggedright\arraybackslash}p{0.180\textwidth}|>{\raggedright\arraybackslash}p{0.585\textwidth}|}
\hline
Sprint & Duration & Main Activities \\ \hline
\endhead
Sprint 1 & Week 1–3 & Requirement analysis, system architecture, database design, and project setup. \\ \hline
Sprint 2 & Week 4–5 & Authentication, user profiles, and university management. \\ \hline
Sprint 3 & Week 6–8 & Forums, posts, comments, interactions, and real-time private/\allowbreak{}group messaging and notifications. \\ \hline
Sprint 4 & Week 9–11 & Document sharing, file search, meeting integration. \\ \hline
Sprint 6 & Week 12 & Integration testing, performance testing, bug fixing, and deployment. \\ \hline
\end{longtable}
\endgroup
\par\medskip
\hypertarget{sec:re-estimation-situation}{}
\addcontentsline{toc}{subsubsection}{Re-estimation Situation}
\subsubsection*{Re-estimation Situation}
\par\medskip
Cost and schedule estimates should be reviewed at the end of each sprint and whenever significant changes occur.
\par\medskip
Re-estimation is necessary in the following situations:
\par\medskip
\begin{itemize}
\item \textbf{Scope changes}: New features or changes to existing requirements increase development effort.
\par
\item \textbf{Technical difficulties}: Unexpected problems with WebSocket, database performance, file sharing, or system integration delay implementation.
\par
\item \textbf{Infrastructure limitations}: Existing server or storage cannot satisfy performance or capacity requirements, requiring additional resources.
\par
\item \textbf{Schedule deviations}: Actual progress differs significantly from the planned milestones, such as a delay exceeding one week.
\par
\item \textbf{External dependencies}: Third-party services introduce additional fees, usage limits, or compatibility issues.
\par
\end{itemize}
During re-estimation, we will compare actual progress and expenses against the initial plan, estimate the remaining work, and adjust the schedule or budget accordingly.
\par\medskip
\hypertarget{sec:42-project-plan}{}
\addcontentsline{toc}{subsection}{4.2 Project Plan}
\subsection*{4.2 Project Plan}
\par\medskip
\hypertarget{sec:421-phase-plan}{}
\addcontentsline{toc}{subsubsection}{4.2.1 Phase Plan}
\subsubsection*{4.2.1 Phase Plan}
\par\medskip
The project is divided into the following major work packages:
\par\medskip
\begin{itemize}
\item \textbf{Project Planning}: Requirement analysis, scope definition, project setup, and architecture design.
\par
\item \textbf{User Management}: Registration, authentication, user profiles, and university management.
\par
\item \textbf{Forum and Social Features}: Forums, posts, comments, interactions, friend requests, and friend management.
\par
\item \textbf{Real-time Communication}: Private/\allowbreak{}group messaging, file sharing, and notifications.
\par
\item \textbf{Document Management}: Document upload/\allowbreak{}download, sharing, public library, and file search.
\par
\item \textbf{Meeting and Administration}: Meeting integration, staff management, roles, and content moderation.
\par
\item \textbf{Testing and Deployment}: Integration testing, performance testing, bug fixing, deployment, and documentation.
\par
\end{itemize}
The project is planned for \textbf{12 weeks}, divided into five sprints.
\par\medskip
\begingroup
\scriptsize
\renewcommand{\baselinestretch}{1.0}\selectfont
\renewcommand{\arraystretch}{1.15}
\setlength{\tabcolsep}{4pt}
\begin{longtable}{|>{\raggedright\arraybackslash}p{0.330\textwidth}|>{\raggedright\arraybackslash}p{0.111\textwidth}|>{\raggedright\arraybackslash}p{0.111\textwidth}|>{\raggedright\arraybackslash}p{0.111\textwidth}|>{\raggedright\arraybackslash}p{0.125\textwidth}|>{\raggedright\arraybackslash}p{0.097\textwidth}|}
\hline
Activities & W1–3 & W4–5 & W6–8 & W9–11 & W12 \\ \hline
\endhead
Requirement Analysis \& Architecture & ■ &  &  &  &  \\ \hline
Database Design \& Project Setup & ■ &  &  &  &  \\ \hline
Authentication \& User Profiles &  & ■ &  &  &  \\ \hline
University Management &  & ■ &  &  &  \\ \hline
Forums, Posts \& Comments &  &  & ■ &  &  \\ \hline
Real-time Messaging \& Notifications &  &  & ■ &  &  \\ \hline
Document Sharing \& File Search &  &  &  & ■ &  \\ \hline
Meeting Integration &  &  &  & ■ &  \\ \hline
Integration Testing &  &  &  &  & ■ \\ \hline
Performance Testing &  &  &  &  & ■ \\ \hline
Bug Fixing \& Deployment &  &  &  &  & ■ \\ \hline
\end{longtable}
\endgroup
\par\medskip
After these sprints, the project must achieve the major milestones below with following criteria:
\par\medskip
\begin{itemize}
\item Project foundation completion: Requirements, architecture, database design, and development environment are ready.
\par
\item User management completion: Users can register, log in, manage profiles, and access university-related features.
\par
\item Core platform completion: Forums, social interactions, real-time messaging, and notifications work correctly.
\par
\item Main features completion: Document sharing, file search, and meeting integration are functional.
\par
\end{itemize}
Corresponding with each milestones, the following demo must be released:
\par\medskip
\begingroup
\scriptsize
\renewcommand{\baselinestretch}{1.0}\selectfont
\renewcommand{\arraystretch}{1.15}
\setlength{\tabcolsep}{4pt}
\begin{longtable}{|>{\raggedright\arraybackslash}p{0.234\textwidth}|>{\raggedright\arraybackslash}p{0.207\textwidth}|>{\raggedright\arraybackslash}p{0.494\textwidth}|}
\hline
Release & Target & Scope and Purpose \\ \hline
\endhead
User Management & End of Week 5 & Demonstrate authentication, user profiles, and university management. \\ \hline
Core Platform & End of Week 7 & Demonstrate forums, social interactions, real-time messaging, and notifications. \\ \hline
Feature Complete & End of Week 9 & Demonstrate document management, file search, and meeting integration. \\ \hline
Final Release: MVP & End of Week 12 & Deploy the application and demonstrate the complete user workflow. \\ \hline
\end{longtable}
\endgroup
\par\medskip
Each demo will be followed by feedback collection and backlog updates. The final release is accepted when all agreed MVP requirements are implemented, critical defects are resolved, and the application passes the required tests.
\par\medskip
\hypertarget{sec:422-iteration-objectives}{}
\addcontentsline{toc}{subsubsection}{4.2.2 Iteration Objectives}
\subsubsection*{4.2.2 Iteration Objectives}
\par\medskip
Each iteration focuses on delivering a specific set of functionalities and achieving measurable development objectives.
\par\medskip
\begingroup
\scriptsize
\renewcommand{\baselinestretch}{1.0}\selectfont
\renewcommand{\arraystretch}{1.15}
\setlength{\tabcolsep}{4pt}
\begin{longtable}{|>{\raggedright\arraybackslash}p{0.122\textwidth}|>{\raggedright\arraybackslash}p{0.122\textwidth}|>{\raggedright\arraybackslash}p{0.691\textwidth}|}
\hline
Iteration & Duration & Objectives \\ \hline
\endhead
\textbf{Sprint 1} & Week 1–3 & - Analyze and finalize project requirements.\newline{}- Define the system architecture and technology stack.\newline{}- Design the database schema and main data relationships.\newline{}- Set up the development environment, project structure, and version control workflow. \\ \hline
\textbf{Sprint 2} & Week 4–5 & - Implement user registration, login, and authentication.\newline{}- Develop user profile management.\newline{}- Implement university management and associate users with their universities.\newline{}- Establish basic role-based access control. \\ \hline
\textbf{Sprint 3} & Week 6–8 & - Implement forum creation and management.\newline{}- Develop post creation, comments, and user interactions.\newline{}- Implement friend requests and friend management.\newline{}- Develop real-time private and group messaging using WebSocket.\newline{}- Implement basic message notifications. \\ \hline
\textbf{Sprint 4} & Week 9–11 & - Implement document upload and download.\newline{}- Enable document sharing in forums and conversations.\newline{}- Develop the public document library and file search functionality.\newline{}- Integrate external meeting services through meeting links. \\ \hline
\textbf{Sprint 5} & Week 12 & - Integrate all system modules and verify end-to-end workflows.\newline{}- Perform functional, integration, and performance testing.\newline{}- Evaluate system behavior under expected concurrent workloads.\newline{}- Fix critical bugs and improve system stability.\newline{}- Deploy the MVP and prepare technical documentation. \\ \hline
\end{longtable}
\endgroup
\par\medskip
At the end of the project, the team expects to deliver a functional MVP that supports university-based communities, social interactions, real-time communication, document sharing, and online meeting integration. The application will be deployed and validated through the agreed testing activities.
\par\medskip
\hypertarget{sec:423-releases}{}
\addcontentsline{toc}{subsubsection}{4.2.3 Releases}
\subsubsection*{4.2.3 Releases}
\par\medskip
The project will follow an incremental release approach. Each release introduces additional functionality and allows stakeholders to review the system before the final deployment.
\par\medskip
\begingroup
\scriptsize
\renewcommand{\baselinestretch}{1.0}\selectfont
\renewcommand{\arraystretch}{1.15}
\setlength{\tabcolsep}{4pt}
\begin{longtable}{|>{\raggedright\arraybackslash}p{0.118\textwidth}|>{\raggedright\arraybackslash}p{0.133\textwidth}|>{\raggedright\arraybackslash}p{0.201\textwidth}|>{\raggedright\arraybackslash}p{0.466\textwidth}|}
\hline
Release & Target & Release Type & Description \\ \hline
\endhead
\textbf{Release 0.1} & End of Week 5 & Alpha & Introduces user registration, authentication, profile management, and university management. \\ \hline
\textbf{Release 0.2} & End of Week 7 & Alpha & Adds forum functionality, social interactions, real-time private/\allowbreak{}group messaging, and notifications. \\ \hline
\textbf{Release 0.3} & End of Week 9 & Beta & Includes document management, file sharing, file search, and external meeting integration. The main features are available for integrated testing and stakeholder feedback. \\ \hline
\textbf{Release 1.0} & End of Week 12 & Initial Production Release (MVP) & Provides the complete agreed MVP after integration testing, performance testing, and critical bug fixing. The application is deployed for initial use. \\ \hline
\end{longtable}
\endgroup
\par\medskip
Each release will be reviewed through a demonstration. Feedback and identified issues will be recorded and prioritized for subsequent iterations.
\par\medskip
\hypertarget{sec:424-project-schedule}{}
\addcontentsline{toc}{subsubsection}{4.2.4 Project Schedule}
\subsubsection*{4.2.4 Project Schedule}
\par\medskip
The project is scheduled to be completed within \textbf{12 weeks}, consisting of five development iterations. The following table summarizes the target completion time for each iteration, release, demonstration, and major milestone.
\par\medskip
\begingroup
\scriptsize
\renewcommand{\baselinestretch}{1.0}\selectfont
\renewcommand{\arraystretch}{1.15}
\setlength{\tabcolsep}{4pt}
\begin{longtable}{|>{\raggedright\arraybackslash}p{0.139\textwidth}|>{\raggedright\arraybackslash}p{0.326\textwidth}|>{\raggedright\arraybackslash}p{0.278\textwidth}|>{\raggedright\arraybackslash}p{0.174\textwidth}|}
\hline
Week & Activity /\allowbreak{} Milestone & Deliverable /\allowbreak{} Release & Event Type \\ \hline
\endhead
Week 1–3 & Requirement analysis, architecture design, database design, and project setup & System foundation & Sprint 1 Completion \\ \hline
End of Week 3 & M1: Project Foundation Completed & Approved requirements, architecture, and database schema & Milestone + Demo 1 \\ \hline
Week 4–5 & Authentication, user profiles, and university management & Release 0.2 – Alpha & Sprint 2 Completion \\ \hline
End of Week 5 & M2: User Management Completed & Functional user and university management & Milestone + Demo 2 \\ \hline
Week 6–8 & Forums, social interactions, real-time messaging, and notifications & Release 0.3 – Alpha & Sprint 3 Completion \\ \hline
End of Week 7 & M3: Core Platform Completed & Functional forum and communication features & Milestone + Demo 3 \\ \hline
Week 9-11 & Document sharing, file search, and meeting integration & Release 0.4 – Beta & Sprint 4 Completion \\ \hline
End of Week 9 & M4: Main Features Completed & All planned core features implemented & Milestone + Demo 4 \\ \hline
Week 12 & Integration testing, performance testing, bug fixing, and deployment & Release 1.0 – MVP & Sprint 5 \\ \hline
End of Week 12 & M5: MVP Released & Tested and deployed application & Final Release + Demo 5 \\ \hline
\end{longtable}
\endgroup
\par\medskip
Progress will be reviewed at the end of each iteration. Any significant deviation from the planned milestones will trigger a schedule review and necessary adjustments.
\par\medskip
\hypertarget{sec:425-project-resourcing}{}
\addcontentsline{toc}{subsubsection}{4.2.5 Project Resourcing}
\subsubsection*{4.2.5 Project Resourcing}
\par\medskip
The following table describes the required roles, estimated staffing, and necessary skills:
\par\medskip
\begingroup
\scriptsize
\renewcommand{\baselinestretch}{1.0}\selectfont
\renewcommand{\arraystretch}{1.15}
\setlength{\tabcolsep}{4pt}
\begin{longtable}{|>{\raggedright\arraybackslash}p{0.260\textwidth}|>{\raggedright\arraybackslash}p{0.190\textwidth}|>{\raggedright\arraybackslash}p{0.486\textwidth}|}
\hline
Role & Number of Staff & Required Skills and Responsibilities \\ \hline
\endhead
Project Manager & 1 & Project planning, progress tracking, risk management, documentation, and team coordination. \\ \hline
Backend Developer & 1–2 & NestJS, TypeScript, REST APIs, authentication, WebSocket, and backend architecture. \\ \hline
Frontend Developer & 1–2 & React, TypeScript, Vite, UI development, state management, and API integration. \\ \hline
Database Engineer & 1 & PostgreSQL, relational database design, SQL, query optimization, indexing, and data management. \\ \hline
AI Engineer & 1 & Python, machine learning fundamentals, model integration, AI APIs, and basic model evaluation. \\ \hline
QA /\allowbreak{} Testing Engineer & 1 & Test case design, functional testing, integration testing, bug reporting, and performance testing. \\ \hline
DevOps /\allowbreak{} Deployment Engineer & 1 & Linux server administration, environment configuration, deployment, and basic system monitoring. \\ \hline
\end{longtable}
\endgroup
\par\medskip
These roles represent the required areas of responsibility rather than separate full-time positions. Depending on the actual team size, members may take on multiple roles.
\par\medskip
The role and responsibilities of team members will be assigned and distributed following these steps:
\par\medskip
\begin{enumerate}
\item Identify each member's existing technical knowledge, practical experience, and areas of interest.
\par
\item Evaluate their suitability for the required roles based on relevant skills and previous projects.
\par
\item Assign roles according to individual strengths while ensuring that all required technical areas are covered.
\par
\item Encourage knowledge sharing and collaboration to reduce dependency on individual members.
\par
\item Review team capacity and role suitability throughout the project, adjusting responsibilities when necessary.
\par
\end{enumerate}
Team members may need additional training to ensure consistent development practices and sufficient technical knowledge.
\par\medskip
\begingroup
\scriptsize
\renewcommand{\baselinestretch}{1.0}\selectfont
\renewcommand{\arraystretch}{1.15}
\setlength{\tabcolsep}{4pt}
\begin{longtable}{|>{\raggedright\arraybackslash}p{0.417\textwidth}|>{\raggedright\arraybackslash}p{0.297\textwidth}|>{\raggedright\arraybackslash}p{0.220\textwidth}|}
\hline
Training Topic & Target Participants & Target Completion \\ \hline
\endhead
Git workflow, branching strategy, and Pull Requests & All team members & Week 1 \\ \hline
NestJS, React, TypeScript, and project architecture & Backend and frontend developers & Week 2 \\ \hline
PostgreSQL, database design, indexing, and query optimization & Database and backend developers & Week 3 \\ \hline
Machine learning fundamentals and AI model/\allowbreak{}API integration & AI Engineer & Week 4 \\ \hline
WebSocket and real-time communication & Backend and frontend developers & Week 5 \\ \hline
File upload, storage, and access control & Backend and database developers & Week 7 \\ \hline
Functional, integration, and performance testing & All team members, especially QA & Week 9 \\ \hline
Linux deployment and server configuration & DevOps /\allowbreak{} Deployment Engineer & Week 10 \\ \hline
\end{longtable}
\endgroup
\par\medskip
Training will primarily use official documentation, online tutorials, and practical exercises. Team members are expected to complete the relevant training before starting tasks that depend on those skills.
\par\medskip
\hypertarget{sec:43-project-monitoring-and-control}{}
\addcontentsline{toc}{subsection}{4.3 Project Monitoring and Control}
\subsection*{4.3 Project Monitoring and Control}
\par\medskip
\hypertarget{sec:431-requirements-management}{}
\addcontentsline{toc}{subsubsection}{4.3.1 Requirements Management}
\subsubsection*{4.3.1 Requirements Management}
\par\medskip
System requirements are documented in the Vision Document. Progress is reviewed at the end of each sprint based on the number of requirements implemented and verified. Any proposed changes, additions, or removals must be submitted as Change Requests and reviewed through the Configuration Management process. Their impact on project scope and schedule will be evaluated before approval. Approved changes will be reflected in the Vision Document and project backlog to ensure that all requirements remain consistent and traceable.
\par\medskip
\hypertarget{sec:432-reporting-and-measurement}{}
\addcontentsline{toc}{subsubsection}{4.3.2 Reporting and Measurement}
\subsubsection*{4.3.2 Reporting and Measurement}
\par\medskip
The team will generate weekly progress reports to track completed tasks, ongoing activities, issues, and upcoming work. At the end of each sprint, a review report will be shared with team members and stakeholders to summarize deliverables, feedback, and project progress. A final project report will be prepared at the end of Week 12.
\par\medskip
The following metrics will be collected to monitor project performance:
\par\medskip
\begingroup
\scriptsize
\renewcommand{\baselinestretch}{1.0}\selectfont
\renewcommand{\arraystretch}{1.15}
\setlength{\tabcolsep}{4pt}
\begin{longtable}{|>{\raggedright\arraybackslash}p{0.381\textwidth}|>{\raggedright\arraybackslash}p{0.571\textwidth}|}
\hline
Metric & Purpose \\ \hline
\endhead
Completed Tasks & Measure development progress. \\ \hline
Requirements Implemented & Track progress against the project scope. \\ \hline
Open and Closed Defects & Monitor software quality and issue resolution. \\ \hline
Test Cases Passed & Evaluate system correctness and readiness. \\ \hline
Schedule Variance & Identify delays compared to the planned schedule. \\ \hline
Actual Expenses & Monitor project spending against the estimated budget. \\ \hline
\end{longtable}
\endgroup
\par\medskip
Reports will be maintained in the shared project workspace and distributed to relevant team members and stakeholders.
\par\medskip
\hypertarget{sec:433-risk-management}{}
\addcontentsline{toc}{subsubsection}{4.3.3 Risk Management}
\subsubsection*{4.3.3 Risk Management}
\par\medskip
Project risks will be identified at the beginning of the project and reviewed at the end of each sprint. Each risk will be assessed based on its likelihood and potential impact, then prioritized as High, Medium, or Low. The Project Manager will maintain the risk register and monitor mitigation actions throughout the project.
\par\medskip
\begingroup
\scriptsize
\renewcommand{\baselinestretch}{1.0}\selectfont
\renewcommand{\arraystretch}{1.15}
\setlength{\tabcolsep}{4pt}
\begin{longtable}{|>{\raggedright\arraybackslash}p{0.135\textwidth}|>{\raggedright\arraybackslash}p{0.414\textwidth}|>{\raggedright\arraybackslash}p{0.386\textwidth}|}
\hline
Risk Ranking & Risk Description and Impact & Mitigation Strategy and/\allowbreak{}or Contingency Plan \\ \hline
\endhead
High & \textbf{Schedule delays:} Complex features may take longer than expected, affecting the final release. & Prioritize core features, monitor weekly progress, and defer non-essential features if necessary. \\ \hline
High & \textbf{Performance issues:} The system may not meet the targets of 2,000 concurrent users and 10,000 requests per second. & Conduct load testing, optimize database queries and APIs, and upgrade infrastructure if necessary. \\ \hline
High & \textbf{Technical integration issues:} Backend, frontend, database, and AI components may not integrate correctly. & Define interfaces early and perform integration testing throughout development. \\ \hline
Medium & \textbf{Security vulnerabilities:} Improper authentication or access control may expose user data. & Apply input validation, authentication, authorization, and security testing. \\ \hline
Medium & \textbf{Data loss:} Database or file storage failures may result in lost user data. & Perform regular backups and test data recovery procedures. \\ \hline
Medium & \textbf{Requirement changes:} New requirements may increase workload and delay development. & Review Change Requests and evaluate their impact before approval. \\ \hline
Low & \textbf{Team coordination issues:} Poor communication may reduce development efficiency. & Assign clear responsibilities and hold regular progress meetings. \\ \hline
\end{longtable}
\endgroup
\par\medskip
\hypertarget{sec:project-close-out}{}
\addcontentsline{toc}{section}{Project Close-out}
\section*{Project Close-out}
\par\medskip
The project will be formally closed after the MVP has been deployed and the final deliverables have been reviewed and accepted.
\par\medskip
The following activities will be performed to ensure an orderly project completion:
\par\medskip
\begin{enumerate}
\item \textbf{Final Acceptance:} Verify that the completed system meets the agreed MVP requirements. Conduct a final demonstration and obtain confirmation from relevant stakeholders.
\par
\item \textbf{Documentation and Archiving:} Complete technical documentation, user guides, database designs, API documentation, and deployment instructions. Archive the source code, project documents, test results, and reports in the shared project repository.
\par
\item \textbf{Final Testing and Reporting:} Summarize testing results, completed features, unresolved issues, and any deviations from the original plan. Prepare a final project report covering project outcomes, schedule performance, and actual expenses.
\par
\item \textbf{Post-mortem Review:} Hold a final meeting to evaluate what went well, what challenges occurred, and what could be improved in future projects. Record lessons learned for future development.
\par
\item \textbf{Team Reassignment and Resource Release:} Since this is a student project, team members will be released from their assigned project responsibilities
\par
\end{enumerate}
\hypertarget{sec:434-configuration-management}{}
\addcontentsline{toc}{subsubsection}{4.3.4 Configuration Management}
\subsubsection*{4.3.4 Configuration Management}
\par\medskip
Git will be used for version control, while a shared issue-tracking system will be used to manage problems and change requests. All reported issues and proposed changes must be reviewed and approved before implementation. Approved changes will be developed in separate branches, tested, and merged through Pull Requests.
\par\medskip
Project artifacts, including source code, database migrations, test scripts, documentation, and configuration files, will be organized using consistent naming conventions. Git tags will identify important versions and releases (e.g., \texttt{v0.1.0}, \texttt{v1.0.0}). Dependencies will be tracked using package manifests and lock files.
\par\medskip
Source code and documentation will be stored in the shared Git repository and cloud storage. Important project files and database data will be backed up periodically. In case of data loss, the team will restore files from the latest stable Git version or available backups. Final source code, documentation, test results, and release artifacts will be archived after project completion for future maintenance.
\par\medskip

\end{document}

